\section{Сообщения оператору}

\subsection{Сообщения, выдаваемые в ходе выполнения программы}

В ходе выполнения программа может выдавать оператору справочные, диагностические
и ошибочные сообщения. Сообщения выводятся в стандартный поток вывода или в
стандартный поток ошибок в зависимости от характера операции.

К справочным сообщениям относятся вывод списка таблиц, списка форматов и полной
справки по параметрам запуска. К сообщениям об ошибках относятся уведомления о
неверных аргументах командной строки, недопустимых значениях параметров,
неизвестных таблицах и неподдерживаемых форматах или режимах выполнения.

\subsection{Действия оператора}

При получении сообщения об ошибке оператор должен:
\begin{enumerate}
    \item проверить корректность имени таблицы, формата и переданных параметров;
    \item проверить допустимость числовых значений параметров;
    \item убедиться, что выбранный формат поддерживается текущей сборкой;
    \item повторить запуск программы после исправления параметров;
    \item при необходимости использовать \texttt{--help}, \texttt{--list-tables} и
    \texttt{--list-formats} для уточнения допустимых значений.
\end{enumerate}

\subsection{Таблица сообщений}

\begin{longtable}{|p{0.28\textwidth}|p{0.34\textwidth}|p{0.28\textwidth}|}
\hline
\textbf{Сообщение} & \textbf{Описание} & \textbf{Действия оператора} \\
\hline
\endfirsthead
\hline
\textbf{Сообщение} & \textbf{Описание} & \textbf{Действия оператора} \\
\hline
\endhead
\texttt{Invalid arguments: ...} &
Ошибка разбора аргументов командной строки. &
Проверить синтаксис команды и повторить запуск. \\
\hline
\texttt{Scale factor must be > 0} &
Задано недопустимое значение масштаба генерации. &
Указать положительное значение параметра \texttt{--scale-factor}. \\
\hline
\texttt{Batch rows must be > 0} &
Задано недопустимое число строк в батче. &
Указать положительное значение параметра \texttt{--batch-rows}. \\
\hline
\texttt{Unknown table: ...} &
Указано имя таблицы, отсутствующее в поддерживаемом наборе. &
Проверить имя таблицы или предварительно выполнить `--list-tables'. \\
\hline
\texttt{Invalid writer options: ...} &
Переданы недопустимые параметры выбранного формата записи. &
Проверить параметры формата и повторить запуск. \\
\hline
\texttt{Invalid write scheduler options: ...} &
Переданы недопустимые параметры планировщика записи. &
Проверить значения \texttt{--write-strategy}, \texttt{--write-workers} и
\texttt{--write-queue-capacity}. \\
\hline
\texttt{Unsupported output format: ...} &
Указан формат, не поддерживаемый текущей сборкой программы. &
Проверить вывод \texttt{--list-formats} и выбрать допустимый формат. \\
\hline
\texttt{Unsupported write-strategy: ...} &
Указана неподдерживаемая стратегия записи. &
Использовать одно из допустимых значений: \texttt{auto},
\texttt{parallel-files}, \texttt{single-file-ordered}. \\
\hline
\texttt{Unsupported verbosity: ...} &
Указан неподдерживаемый режим вывода. &
Использовать \texttt{quiet} или \texttt{verbose}. \\
\hline
\texttt{Table task failed: ...} &
Ошибка выполнения одной из задач генерации таблицы. &
Проверить параметры запуска, доступность ресурсов и повторить запуск. \\
\hline
\end{longtable}
